\section{Introduction}
Mean calibration leaves upper-tail agreement unexamined when a cheap instrument accompanies a reference instrument.
The air-monitoring study uses a linear calibration trained on historical collocations.
In the subsequent audit period, its mean difference is $-0.191$ micrograms per cubic metre near 30\% relative humidity (RH).
The calibrated sensor assigns 6.068\% of local observations above concentration 15, against 9.307\% for the reference instrument.
Distributional verification asks which discrepancies can be established from a limited reference sample.

Covariates $X$ and a surrogate $S$ are widely observed, while a reference outcome $Y$ is selectively verified.
The target is the conditional cumulative distribution function (CDF) of $Y$ at a chosen $x_0$.
Augmented inverse-probability weighting remains valid when the fitted outcome regression is inaccurate.
A common borrowing path joins surrogate-free Anchor to a surrogate-assisted estimator.

For a simultaneous band, average variance is the wrong criterion.
The band's critical value depends on the joint covariance across thresholds, including how the correction redistributes variance.
Proposition~\ref{prop:separation} gives a valid two-threshold path whose population trace-minimising weight reduces total variance and increases the simultaneous critical value.
Estimating a band-optimal weight adds covariance uncertainty, since a small apparent reduction may reflect sampling noise.

Surrogate-outcome and two-phase sampling establish the observed-data representation for selectively measured outcomes \citep{Pepe1992,RobinsRotnitzkyZhao1994,TaoZengLin2017,Tsiatis2006}.
Adaptive augmentation and semi-supervised inference optimise auxiliary-information use for parameter estimation \citep{CaoTsiatisDavidian2009,ChakraborttyCai2018,ZhangBrownCai2019,Song2024}.
Prediction-powered inference places this construction in a broad prediction-assisted framework \citep{AngelopoulosEtAl2023,AngelopoulosDuchiZrnic2023,ZrnicCandes2024}.
Related work documents the finite-reference cost of estimating borrowing weights \citep{mani2025nofreelunch,Eyre2025,KallusMao2025}.
These procedures choose auxiliary information for parameter estimation; the selection criterion for a threshold-indexed local distribution has remained open.

\citet{SuiEtAl2026} study prediction-powered conditional functionals, and \citet{TestaEtAl2025} allow selective labels with decreasing overlap.
For transported outcomes, \citet{Wang2026} develops a framework for surrogate-assisted distributional inference.
Local distribution estimation with fully observed outcomes has an established theory \citep{CattaneoJanssonMa2024}.
Generation-powered and multi-source extensions develop functional borrowing operators and confidence-region volume \citep{zhang2026generation,li2026multisource}.
Anderson's covariance-dominating operator improves Gaussian probabilities of centred convex symmetric sets \citep{Anderson1955}, and a fitted path may lack this ordering.

We develop \emph{paired gain certification} through a common covariance comparison.
The estimated object is the difference between Anchor and candidate simultaneous critical values, whose paired uncertainty vanishes at zero borrowing.
Exploiting this cancellation yields an oracle bound proportional to the beneficial borrowing weight.
An implementable multiplier calibration with independent numerical guards supports simultaneous inference after selection with nonunique optimal weights.
A binary local experiment gives the reference-information scale for certifying weak improvements.
The general learning-and-testing principle is established \citep{angelopoulos2021learn}; our contribution is the paired critical-value geometry, its local information scaling, and its use in distributional inference.

Section~2 defines the estimator and algorithm, and Section~3 states the certification, inference and information results.
Section~4 compares borrowing criteria and evaluates their extensions; Section~5 analyses dated air-monitoring collocations before the concluding discussion.
